% Full source/proof audits passed by root and frobenius.
\paragraph{E.3t (Proved: the full morphic obstruction to Albanese equivalence).}
Let $X/\mathbb C$ be smooth projective of dimension $n$, let
$0\le i\le n$, and put $r=n-i$. We distinguish the full
morphic Jacobian
\[
 J^{\mathrm{mor}}_r(X)
   =J\bigl(L_rH_{2r+1}(X,\mathbb Z)\bigr)
\]
from the finite-dimensional Walker Jacobian $J_W^i(X)$.
Here $L_rH_m$ is Lawson homology with the inductive mixed
Hodge structure of \citet[Sections 3--5]{cyWalker2007Morphic};
the convention already incorporates the twist making the
relevant pure part have weight $-1$. There are canonical
surjections
\[
 P^i(X)\longrightarrow J^{\mathrm{mor}}_r(X)
       \longrightarrow J_W^i(X)(\mathbb C)
       \longrightarrow J_a^i(X)(\mathbb C).
\]
The last map is an isogeny of abelian varieties. No such
finite-dimensionality assertion applies to the first
morphic target. In particular
\begin{equation}\label{eq:E3-N-killed-morphic}
        N^i(X)\subseteq\ker\Phi^{\mathrm{mor}}_r.
\end{equation}

More precisely, for an abelian variety $A/\mathbb C$ and
$Z\in\operatorname{CH}^i(A\times X)$, let
\[
 \lambda_Z:H_1(A,\mathbb Q)
           \longrightarrow L_rH_{2r+1}(X,\mathbb Q)
\]
be the Lawson cylinder map induced by this correspondence.
Then the composite of $u_Z$ with $P^i(X)\to
J^{\mathrm{mor}}_r(X)$ is zero if and only if
$\lambda_Z=0$. Consequently
\begin{equation}\label{eq:E3-Lawson-necessary}
 u_Z=0\ \Longrightarrow\ \lambda_Z=0
       \ \Longrightarrow\ a({}^tZ)=0.
\end{equation}
The same first implication holds if $u_Z\otimes\mathbb Q=0$.
The next paragraph shows that the second implication cannot
be reversed even for an elliptic parameter.

\begin{proof}
Example 5.5 of \citet[p.~928]{cyWalker2007Morphic} identifies
the morphic Abel--Jacobi map of zero-cycles on every smooth
projective $T$ with its ordinary Albanese map.
Theorem 5.7(2),(4), printed p.~930, says that the morphic
map factors through rational equivalence and is functorial
for all correspondences between smooth projective varieties
with nonnegative dimension shift. A generator of $N^i(X)$
is $W_*\nu$, where
\[
 W\in\operatorname{CH}^i(T\times X),\qquad
 \nu\in\operatorname{CH}_0(T)^0,\qquad
 \operatorname{alb}_T(\nu)=0.
\]
Its dimension shift is $r\ge0$, so the theorem applies and
gives $\Phi^{\mathrm{mor}}_r(W_*\nu)=0$. This proves
\eqref{eq:E3-N-killed-morphic}. Theorem 5.7(1) gives
surjectivity from algebraically trivial cycles to the full
morphic Jacobian. The factorization through the finite
Walker Jacobian is the displayed factorization immediately
before Corollary 5.9, printed p.~932, and that corollary
gives the finite kernel of $J_W^i\to J_a^i$.

We include the point that detects an entire family. Put
$H=L_rH_{2r+1}(X,\mathbb Z)$. It is a countable inductive
limit of finitely generated mixed Hodge structures, so its
underlying group is countable. By Proposition 3.3, printed
pp.~916--917, its Jacobian is
\[
 J(H)=\frac{H_{\mathbb C}}{F^0H_{\mathbb C}+\operatorname{im}H}
\]
in these weights. The correspondence map from $H_1(A)$ is
a morphism of inductive mixed Hodge structures. If its
rationalization is nonzero, its image is a nonzero finite
pure Hodge structure of weight $-1$. Strictness of the
Hodge filtration implies that the induced complex-linear
map
\[
 H_1(A,\mathbb C)/F^0
               \longrightarrow H_{\mathbb C}/F^0
\]
has nonzero image. That image cannot be contained in the
countable image of $H$, or even in its rational span.
Thus the induced Jacobian map has nontorsion image.
Conversely, if the rational map is zero, the integral
map has finite torsion image; the Jacobian of that finite
Hodge structure is zero, so the induced Jacobian map
vanishes. Functoriality and the zero-cycle comparison
identify this Jacobian map with
$a\mapsto\Phi^{\mathrm{mor}}_r(Z_*([a]-[0]))$.
This proves the equivalence and the first implication in
\eqref{eq:E3-Lawson-necessary}, including its rational
version. Composition with the Lawson-to-singular cycle
map gives the ordinary cohomological cylinder. Poincar\'e
duality identifies its transpose with $a({}^tZ)$.
\end{proof}

\paragraph{E.3u (Refuted: the complex converse, with a specified sixfold and elliptic family).}
The implication
\[
                     a({}^tZ)=0\ \Longrightarrow\ u_Z=0
\]
of 14:38--39 is false over $\mathbb C$, even after
rationalizing its conclusion. We give a codimension-four
counterexample on a smooth projective sixfold, with an
elliptic curve as parameter and with an explicitly
specified point whose class modulo $N^4$ is nontorsion.

\emph{The equations and the cycle.}
Let
\[
 Q=\{x_0y_0+x_1y_1+x_2y_2+x_3y_3=0\}
                  \subset\mathbb P^7_{\mathbb Q},
\]
and set
\[
 \Pi_+=\{x_0=x_1=x_2=x_3=0\},\qquad
 \Pi_-=\{y_0=x_1=x_2=x_3=0\},\qquad
 \zeta=[\Pi_+]-[\Pi_-].
\]
Take the degree-$14$ monomials not divisible by $x_0y_0$
as a basis of $H^0(Q,\mathcal O_Q(14))$. Give $x_0^{14}$
coefficient one and every other basis monomial an
independent indeterminate $t_b$. Put
\[
 K=\mathbb Q(t_b),\qquad \Omega=\overline K,
 \qquad F=x_0^{14}+\sum_{b\ne x_0^{14}}t_b b,
 \qquad Y=Q_\Omega\cap\{F=0\}.
\]
This is a specified geometric generic smooth hypersurface,
of dimension five. The intersections with both $\Pi_\pm$
are proper, and we put
\[
 \gamma=[Y\cap\Pi_+]-[Y\cap\Pi_-]
                   \in\operatorname{CH}_2(Y).
\]
Finally set
\[
 E:\ v^2=u^3-u,\qquad o=\text{the point at infinity},
 \qquad X=Y\times E,
\]
and on the ordered product $E_a\times Y\times E_b$ set
\begin{equation}\label{eq:E3-explicit-morphic-family}
 Z=p_Y^*\gamma\cdot p_{ab}^*
                   (\Delta_E-E_a\times\{o\})
             \in\operatorname{CH}^4(E_a\times X).
\end{equation}
For $a\in E(\mathbb C)$ this gives
\[
                 u_Z(a)=\gamma\times([a]-[o])\pmod{N^4(X)}.
\]
All the stated properties hold after every embedding of
the indicated algebraically closed fields into $\mathbb C$.

\emph{The needed strengthened nonvanishing.}
The quadric calculation in C.1i gives
\[
 \deg\zeta^2=-2,
 \qquad h\operatorname{cl}(\zeta)=0,
 \qquad \operatorname{cl}(\zeta)\ne0.
\]
Integral weak Lefschetz identifies $H_4(Y,\mathbb Z)$
with $H_4(Q,\mathbb Z)\simeq\mathbb Z$. The two
surfaces defining $\gamma$ have the same degree, namely
$14$, so
\begin{equation}\label{eq:E3-gamma-lawson}
 \operatorname{cl}(\gamma)=0\text{ in }H_4(Y,\mathbb Z),
 \qquad
 s\gamma\ne0\text{ in }L_1H_4(Y,\mathbb Q).
\end{equation}
Here the second assertion is stronger than Griffiths
nonvanishing, and requires a separate argument.

We use \citet[Theorem 3.4, printed p.~1082, with its proof
on pp.~1095--1097]{cyFriedlander2000}. This is the theorem
behind Walker's Theorem 6.1; the latter's locator ``4.5''
does not refer to the required theorem in the published
Friedlander paper. We spell out why the effective degree
and the geometric generic field are legitimate. For
$S=\mathbb P H^0(Q,\mathcal O_Q(14))$ and every smooth
base $T\to S^\circ$, the Nagel bound verified in C.1i
gives
\begin{equation}\label{eq:E3-effective-Lawson-connectivity}
 H^k(Q\times T,Y_T;\mathbb Q)=0\qquad(k\le10).
\end{equation}
In Nagel's notation we have $c=\mu=5$, $r=0$, and
$m_Q=1$. The two numerical conditions reduce to
$14\ge6$ and $14\ge12$. Thus no unspecified large-degree
threshold is being used.

For completeness, suppose that a nonzero multiple of
$\gamma$ has zero $s$-image. Theorem 3.2 and
Proposition 5.1 of the same Friedlander paper express
such a relation by a projective threefold $W$, allowed
to be singular, a homologically trivial one-cycle
$\delta$ on $W$, and a Chow correspondence from $W$
to the curve space of $Y$. Finite sums are combined
using disjoint unions. These are finite algebraic data.
Spread them on a smooth dominant base $T\to S^\circ$
and shrink so that the proper families and their
resolutions have the required topological strata and
so that the zero homology class of $\delta$ persists.
The finite-type relative Chow construction in
Proposition 5.1 supplies precisely this spread.
We do not assert that coefficients of an arbitrary
complex witness are algebraic over $K$: a smooth
dominant base is sufficient, since
\eqref{eq:E3-effective-Lawson-connectivity} permits
every such base. The smooth locus of that dominant
map is nonempty in characteristic zero.

We recall the decisive cohomology argument, to identify
the only connectivity degrees used. Resolve the
proper family $W$ by $W'$ and resolve
$W'\times_W W'$ by $W''$; shrink $T$ so their
composites to $T$ are smooth. The relative Chow
correspondence has a characteristic class of degree
eight on $W'\times_TY_T$. Connectivity lifts it
uniquely to a class $b$ of degree eight on
$W'\times Q$, and its two pullbacks to
$W''\times Q$ are equal, by the same uniqueness.
The degree-six injectivity from connectivity says
that its action on the lifted cycle $\delta'$ is
the pullback of the nonzero multiple of
$\operatorname{cl}(\zeta)$. On any fibre, proper
homological descent, as used in the proof on p.~1097,
gives
\[
 [\delta']=(p_1)_*\eta-(p_2)_*\eta
                  \quad\text{for some }\eta\in H_2(W'',\mathbb Q),
\]
because its pushforward $[\delta]$ on $W$ is zero.
The projection formula and $p_1^*b=p_2^*b$ make
the action of $b$ on $[\delta']$ zero. This would
give $m\operatorname{cl}(\zeta)=0$ for a nonzero
integer $m$, a contradiction.
The only restriction degrees are eight and six;
relative vanishing through degree nine suffices,
and \eqref{eq:E3-effective-Lawson-connectivity}
gives vanishing through ten. This proves the second
assertion of \eqref{eq:E3-gamma-lawson}.

This argument also explains transfer for the specified
field. A witness over an extension is spread over a
finite-type model of that extension; proper
topological stratification, or its constructible
cohomology formulation, preserves its zero cycle
class after shrinking. The projection to $S$ is
dominant because the coefficients $t_b$ are
algebraically independent. Its smooth locus is
therefore available, and the preceding argument
applies there. It contradicts a witness after any
embedding into $\mathbb C$, without assuming that
an unspecified analytic exceptional locus is
defined over $\overline{\mathbb Q}$.

\emph{The whole cylinder vanishes, while the family survives.}
Equation \eqref{eq:E3-gamma-lawson} implies
$\operatorname{cl}(Z)=0$ integrally by the exterior
product formula, and therefore
\[
 Z_*=0:H_1(E,\mathbb Q)\longrightarrow
                         H_5(X,\mathbb Q).
\]
Equivalently the cohomological cylinder
$H^1(E,\mathbb Q(1))\to H^7(X,\mathbb Q(4))$
is zero, and
$a({}^tZ)=0$ in the manuscript's degree.

On Lawson homology the same family gives
\[
 \lambda_\gamma:H_1(E,\mathbb Q)\longrightarrow
                       L_2H_5(X,\mathbb Q),
              \qquad v\longmapsto\gamma\times v.
\]
This is injective. Indeed choose a class $w$ on the
elliptic curve with nonzero intersection pairing
against a given $v\ne0$. Cap with the pullback of
$w$ and push to $Y$. The product calculation in
\citet[proof of Theorem 6.2, printed pp.~933--934]
{cyWalker2007Morphic} gives
\[
 (p_Y)_*\bigl((\gamma\times v)\cap p_E^*w\bigr)
                      =\langle v,w\rangle\,s\gamma\ne0.
\]
The nondegenerate intersection form on $H_1(E)$
and \eqref{eq:E3-gamma-lawson} prove injectivity.
Strictness of the Hodge filtration now injects the
one-dimensional complex space $H_1(E,\mathbb C)/F^0$
into the Lawson complex quotient used in E.3t.
The integral Lawson group is countable, so the
induced map
\[
 E(\mathbb C)\longrightarrow J^{\mathrm{mor}}_2(X),
           \qquad a\longmapsto
              \Phi^{\mathrm{mor}}_2(\gamma\times([a]-[o]))
\]
has countable kernel and countable inverse image
of the torsion subgroup. In particular it has
nontorsion image. Since it kills $N^4(X)$,
$u_Z\otimes\mathbb Q\ne0$. This proves the
counterexample for the entire family.

We can specify a point witness as well. Adjoin one
further indeterminate $t$ to $\Omega$, choose
$v_t$ with $v_t^2=t^3-t$, put
\[
 \Omega'=\overline{\Omega(t)},\qquad
 e=(t,v_t)\in E(\Omega'),\qquad
 z=\gamma\times([e]-[o])\in\operatorname{CH}^4(X_{\Omega'}).
\]
Then $z$ is nontorsion modulo $N^4$ after every
embedding into $\mathbb C$. Suppose that $mz\in N^4$
for $m\ne0$. An expression by finitely many defining
Albanese-kernel correspondences, including its
rational-equivalence witnesses, spreads over a
smooth dominant finite-type base above an open
of $E_\Omega$. Smooth projective parameter
varieties, their relative Albanese schemes, the
marked zero-cycles and the correspondence cycles
all spread after shrinking. The equation that
each Albanese class is zero is an equation of
sections of its relative abelian scheme and thus
continues to hold. The Chow relation also
continues to hold because its finite rational-
equivalence witnesses have been spread.
Consequently an entire nonempty complex open
subset of $E$ would have morphic image killed by
$m$. This contradicts the countability just
proved. This descent argument permits auxiliary
coefficients transcendental over $\Omega'$; a
dominant parameter base suffices. It proves the
asserted point witness and transfer.

The finite Walker family map is zero, since it
is an algebraic homomorphism from $E$ with zero
map on $H_1$. Thus on this same sixfold
\[
 N^4(X)\subsetneq\ker\psi_W^4
                    =(\mathcal A*\mathcal H)^4(X),
 \qquad
 \ker\bigl(P^4(X)\to J_W^4(X)(\mathbb C)\bigr)
                         \otimes\mathbb Q\ne0.
\]
By E.3o the proposed general equality
$\langle\mathcal H^{*2}\rangle_0=\mathcal A*\mathcal H$
is therefore false. This does not contradict the
proved surface-support cases E.3j, E.3l and E.3q,
nor does it refute the separate possible comparison
between the large one-motivic quotient and the
\emph{full} morphic Jacobian in E.3s. The latter
comparison is an additional question, not needed
to decide the requested converse.\hfill$\square$

The finite certificate
\texttt{code/check\_morphic\_counterexample.py} verifies
the quadric equations and gradients, the two planes,
the self-intersection coefficient and $\zeta^2=-2$,
the $65\,892$ basis monomials, the degree-$14$
connectivity inequalities, the elliptic discriminant
$64$, and the unimodular intersection matrix.
Its output records the exact runtime. The infinite
generic-fibre and Lawson assertions follow from
the proof, not from the bounded computation.
